
\chapter{Selective Provenance Expansion}
\label{ch:spe}

Graph triples give a compact index of past dialogue, but extraction can remove qualifiers, temporal
context and other details a question needs. Restoring the source turn keeps those details at a higher
token cost. We designed and pre-registered \emph{Selective Provenance Expansion} (SPE) to choose
between the two: under a retrieval-context token budget, it decides whether to omit a retrieved fact,
keep it as a triple, or include its source dialogue turn. This chapter describes the two versions of
the selection rule, evaluates the repaired implementation against the corrected baselines, and
examines why the expected trade-off did not appear.

\paragraph{Contribution of this chapter.} SPE is proposed in this thesis. The closest systems also
return from a compact memory unit to its source text (Section~\ref{sec:ch2_graphs}). SEEM does so
for every retrieved event frame~\citep{seem}, and LeanMem does so when an LLM planner judges that a
question needs detail~\citep{leanmem}. SPE differs in where the decision is made. It weighs the
two representations of each retrieved fact against each other under one shared token budget, and
it was evaluated against criteria registered before its record runs. The result is negative. The
chapter's second contribution is the explanation of that result: the rule's per-token scoring,
together with the measured token costs, sets a bar for source text that the tested inputs did not
clear (Section~\ref{sec:spe_why_triples}).

\section{Why source text may matter}
\label{sec:spe_motivation}

A long-term conversational memory has to trade compactness against fidelity:

\begin{enumerate}
  \item \textbf{Compactness}: Longer prompts increase prefill work and KV-cache storage. Within the
        operating range measured in this thesis, they also increased inference energy. In the SPE runs,
        a selected triple cost 17.5 to 19.2 tokens on average, while upgrading it to its source text
        cost about 75. The shorter representation allows more retrieved facts to fit in the context
        window.
  \item \textbf{Fidelity}: A triple does not preserve every detail of its source turn. Qualifiers,
        temporal expressions, speaker intent, and surrounding context may disappear during extraction,
        leaving the reader with less evidence than the original dialogue contained.
\end{enumerate}

The tension is not an artefact of small extractors: \citet{rag_vs_graphrag} find the same pattern
with a 70B reader, while small-model extraction itself lowers graph-reasoning
accuracy~\citep{kelesis2026}.

The current E2B results in Chapter~\ref{ch:graph_mechanisms} show the trade-off: v3 reaches
21.5\%, source-restored v4 reaches 33.0\%, and vector retrieval reaches 36.0\%.
SPE was specified before these reruns.\footnote{At the design freeze, the
earlier v3 and v4 baselines were 22.0\% and 30.0\%; the registered criteria retain
those values.} Because v3 and v4 select the same graph edges, the gain comes from what the reader
is shown.

This tension motivated the core question of this chapter: \emph{Under an explicit retrieval-context
budget, can an algorithm selectively decompress only the facts that require verbatim evidence,
retaining unambiguous facts as compact triples, to achieve high accuracy at a lower resource footprint?}

SPE uses the repaired Gemma~4 E2B and E4B v3 graph caches described in
Section~\ref{sec:graph_topology}; this chapter concerns how to select representations
from those fixed inputs.

SPE is therefore not a separate method to be combined with GraphRAG v4. It runs on the graph and
the source links that v4 uses. v3 shows every selected fact as a triple only, and v4 attaches the
source turns of every selected fact; SPE makes that choice fact by fact. It selects from a larger
candidate pool (up to 40 edges) under a token budget instead of a fixed twelve, and v4 is the
accuracy reference it was required to approach at a lower prompt cost.

\section{SPE design}
\label{sec:spe_formulation}

SPE was designed as a small change to the existing graph retrieval path, not as a learned routing
system. The existing pipeline already stored three useful pieces of information: the similarity between
the query and each retrieved dialogue turn, the triples extracted from that turn, and the link from each
triple back to its source. We used those values to decide how much of each fact to place in the prompt.

Candidate construction follows the implementation. Dense retrieval first selects the top-$k$ dialogue
turns. SPE then collects graph edges whose provenance overlaps those turns, orders the edges by the
similarity of their best matching turn, and keeps at most 40 candidates. Each candidate can be omitted,
shown as a triple, or shown as a triple together with the retrieved source turn. Source quotes shared by
several triples are included only once, so the cost of a choice depends on what has already been selected.

Algorithm~\ref{alg:spe} summarizes the resulting greedy selection rule.

\begin{algorithm}[htbp]
\caption{Selective provenance expansion using a greedy density rule}
\label{alg:spe}
\small
\begin{algorithmic}[1]
\REQUIRE Query $q$, candidate facts $\mathcal{E}$, retrieved source turns, token budget $B$
\STATE Mark every candidate as omitted and set the spent budget to zero.
\WHILE{some representation change fits within $B$ and has $\Delta U > 0$}
  \FORALL{valid changes: omit$\rightarrow$triple, omit$\rightarrow$expanded, triple$\rightarrow$expanded}
    \STATE Compute the additional tokens $\Delta C$, excluding source quotes already present.
    \STATE Compute the additional utility $\Delta U$ supplied by the change.
  \ENDFOR
  \STATE Among the changes that fit and have $\Delta U > 0$, choose the one with the largest $\Delta U/\Delta C$.
  \STATE Apply the change and update the spent tokens, included source turns, and covered terms.
\ENDWHILE
\STATE \textbf{return} the selected triples followed by the selected source quotes.
\end{algorithmic}
\end{algorithm}

The two SPE versions differ only in how they define utility.

\subsection{SPE-1: a simple heuristic score}
The first version combined signals already available in the retrieval pipeline. Similarity measured
query relevance, lexical overlap rewarded uncovered query terms, a binary flag favored timestamped
facts for temporal questions, and a redundancy term discouraged repeated evidence:
\begin{equation}
  U_1(e, r, \mathcal{S}, \mathcal{C}) = \text{sim}(q, s_e) + \lambda_1 \cdot \text{Cov}(e, q \setminus \mathcal{C}) + \lambda_2 \cdot \text{Temp}(e, q) - \lambda_4 \cdot \text{Red}(e, \mathcal{S})
\end{equation}
Here, $\mathcal{S}$ is the set of source turns already included and $\mathcal{C}$ is the set of query
terms already covered. The similarity, temporal and redundancy terms were identical for a fact's triple
and expanded forms; only the coverage term could differ, because the expanded form also contains the
source text. This issue is examined below.
The three weights were set to 0.3 before evaluation and were not tuned on the test questions. This was
an intentionally simple starting point: it tested whether readily available metadata was enough to
choose which facts deserved more context.

\subsection{SPE-2: score the two representations separately}
SPE-1 had a basic problem. The triple and expanded versions of a fact shared the
same similarity term even though they contained different text. Apart from lexical coverage, the
selector therefore had little information about whether paying for the source turn was worthwhile.
Before running the SPE-2 cells, we replaced that shared similarity with one for the triple and another
for the source turn:
\begin{align}
  u_{\text{triple}}(e) &= \cos(\mathbf{e}_q, \mathbf{e}_{\text{triple}(e)}) \\
  u_{\text{source}}(e) &= \cos(\mathbf{e}_q, \mathbf{e}_{\text{source}(e)}) \\
  \text{gap}(e) &= u_{\text{source}}(e) - u_{\text{triple}}(e)
\end{align}

The difference between the two scores is a proxy for detail that may have been lost during extraction.
An upgrade from a selected triple to its source turn competes using $\text{gap}(e)/\Delta C$. SPE-2
removes the three heuristic weights, but it retains the candidate cap and token budget. The change was
recorded before the SPE-2 evaluation and no further selection rule was introduced afterward.

\section{Evaluation against the pre-registered criteria}
\label{sec:spe_results}

The registered evaluation used Gemma~4~E2B on 200 LoCoMo questions, with answers judged by a 70B
model. A frozen addendum repeated the comparison with E4B. The retrieval-context budget $B$ covers
retrieved triples and source quotes; total prompt length also includes the instructions, answer format,
and question.

The registered criteria used the baseline values available when the design was frozen:
\begin{itemize}
  \item \textbf{Criterion 1}: reach at least 30.0\% accuracy with at least 25\% fewer total prompt
        tokens than the expand-all baseline;
  \item \textbf{Criterion 2}: exceed 22.0\% accuracy while keeping total prompt length below 600 tokens.
\end{itemize}

The fixed criteria remain unchanged; the comparisons below use the current baselines
under the convention in Section~\ref{sec:retrieval}.

\subsection{The selector repair and why the repaired rule is the one evaluated}
\label{sec:spe_repair}

After the first SPE-2 runs and before the record runs, a code audit found two bookkeeping defects in
the selector. The registered design states that choosing the expanded
form of a fact \emph{supersedes} its triple, and that the greedy step takes the remaining option with
the highest utility per token (Algorithm~\ref{alg:spe}). The code departed from that design in two ways:
\begin{enumerate}
  \item Choosing the expanded form of an already selected fact appended a second copy of the fact
        instead of replacing the triple, so the triple line appeared twice in the prompt.
  \item An upgrade from triple to expanded form was charged only its extra tokens, but it was
        credited with the full utility of the expanded form, $u_{\text{source}}$. The triple's
        utility, $u_{\text{triple}}$, had already been counted when the triple was selected, so
        the upgrade's value was counted partly twice.
\end{enumerate}
The repair made the expanded form replace the triple and credited an upgrade with what it adds,
$\text{gap}(e) = u_{\text{source}} - u_{\text{triple}}$, divided by what it adds in tokens. The
first change implements the supersession that the recorded design states. The second is the
implementation's reading of that design: the design sets the expanded form's utility to
$u_{\text{source}}$ and selects greedily by utility per token, but it does not write out how an
upgrade from an already selected triple is credited. In a greedy
budgeted selection, every step must be scored by the value it adds per token it adds. When an
upgrade is credited with value the prompt already contains, the selector pays for source text for
an accounting reason, not because the source text is worth its cost. The repaired rule can still
choose source text. It does so whenever an upgrade's gap per extra token exceeds what the best
remaining triple offers per token, or once every candidate triple has been taken and budget remains.
That it never did so is therefore a result about these inputs, and Section~\ref{sec:spe_why_triples}
explains it.

The earlier version, which did quote source turns, also did not succeed. At $B=370$ it quoted at
least one source turn in 39 of the 200 E2B prompts, and it reached 26.0\%, against 27.0\% for the
repaired rule; both are below the 30.0\% that Criterion~1 requires.\footnote{At $B=240$ the earlier
version reached 24.0\% for E2B, and 24.5\% and 27.0\% for E4B at the two budgets. These earlier runs
also used embeddings without the task prefixes, so they differ from the record runs in two respects
and are not used elsewhere in this chapter.} Neither version met Criterion~1. Because the embedding
prefixes changed at the same time, the repair's own effect on accuracy cannot be isolated. The repair defines the valid implementation used in this chapter and is not a new SPE variant;
the registered design had already declared SPE-2 the final variant.

\subsection{Registered E2B results}
Table~\ref{tab:spe_results} reports the registered E2B comparison.

\begin{table}[htbp]
  \centering
  \scriptsize
  \setlength{\tabcolsep}{3pt}
  \begin{tabular}{lccccc}
    \toprule
    Configuration & Retrieval control & Total Prompt Tokens & Kept Triples & Source Turns & Accuracy \\
    \midrule
    Full context                       & 4096-token window       & 4096.0 & n/a   & n/a   & 7.5\% \\
    GraphRAG v3 (triples)              & cap: 12 triples         & 519.6  & 12.00 & 0.00  & 21.5\% \\
    GraphRAG v4 (triples + sources)    & cap: 12 triples         & 812.7  & 12.00 & 4.50  & 33.0\% \\
    Vector RAG                         & top-10 dialogue turns   & 824.4  & 0.00  & 10.00 & \textbf{36.0\%} \\
    SPE-2 tight                        & context budget: 240 tokens & 545.7  & 13.81 & 0.00  & 23.0\% \\
    SPE-2 registered                   & context budget: 370 tokens & 737.1  & 20.68 & 0.00  & 27.0\% \\
    SPE-1                              & context budget: 240 tokens & 547.4  & 14.03 & 0.00  & 15.0\% \\
    SPE-1                              & context budget: 370 tokens & 738.3  & 20.85 & 0.00  & 20.5\% \\
    \bottomrule
  \end{tabular}
  \caption{Corrected baseline, repaired SPE-2 and SPE-1 results on LoCoMo ($n=200$, Gemma~4~E2B reader,
           70B judged accuracy). The baselines use fact-count or context-window controls, not
           SPE's explicit retrieval-context budget. Token and retrieval counts are means over 200 questions.
           The SPE-1 rows are exploratory and outside the registered criteria.}
  \label{tab:spe_results}
\end{table}

Criterion 1 failed. At $B=370$, SPE-2 reached 27.0\%, compared with 33.0\% for GraphRAG v4, and
reduced mean prompt length by only 9.3\%. The paired difference did not separate (McNemar $p=0.0576$,
Holm-adjusted $p=0.1152$), but the configuration missed both registered thresholds.

Criterion 2 passed numerically but remained inconclusive. The tighter configuration reached 23.0\%
within the prompt limit, yet its 1.5-point gain over GraphRAG v3 was not statistically distinguishable
($p=0.711$).

SPE-1, run on the same pipeline as an exploratory comparison, did worse at both budgets: 15.0\% at
$B=240$ and 20.5\% at $B=370$. Like SPE-2, it selected no source turns. At the tighter budget it fell
below GraphRAG v3 itself (exact McNemar $p=0.029$, unadjusted), so the simpler scoring rule was not
merely unable to recover the benefit of source text but spent its budget less well than the triple cap
it replaced.

\subsection{E4B addendum and reader--extractor factorization}
\label{sec:scale_generalization}

The frozen scale addendum tested whether E4B could make better use of selectively compressed evidence.
Because E4B initially served as both extractor and reader, this was an end-to-end comparison rather
than a reader-only test. Table~\ref{tab:scale_comparison} places the registered E2B result beside the
E4B addendum.

\begin{table}[htbp]
  \centering
  \scriptsize
  \setlength{\tabcolsep}{3pt}
  \begin{tabular}{lcc}
    \toprule
    Memory Configuration & Gemma 4 E2B (2B) & Gemma 4 E4B (4B) \\
    \midrule
    GraphRAG v3 (triples) & 21.5\% & 27.5\% \\
    SPE-2 repaired ($B=240$) & 23.0\% & 25.5\% \\
    SPE-2 repaired ($B=370$)& 27.0\% & 25.0\% \\
    SPE-1 ($B=240$) & 15.0\% & 19.0\% \\
    SPE-1 ($B=370$) & 20.5\% & 22.5\% \\
    GraphRAG v4 (triples + sources) & 33.0\% & 32.5\% \\
    Vector RAG & \textbf{36.0\%} & \textbf{40.5\%} \\
    \bottomrule
  \end{tabular}
  \caption{End-to-end scale sensitivity between Gemma~4~E2B and E4B on LoCoMo ($n=200$). For
           graph-based configurations, each column initially used the same model for extraction and
           reading, so the cross-column difference is not a reader-only effect. The E4B SPE
           comparisons are exploratory.}
  \label{tab:scale_comparison}
\end{table}

E4B did not rescue SPE. Its two SPE-2 conditions reached 25.5\% and 25.0\%, and its two SPE-1
conditions 19.0\% and 22.5\%, all below GraphRAG v4 at 32.5\% and Vector RAG at 40.5\%. The larger configuration therefore preserved the main result rather
than narrowing the representation gap. Figure~\ref{fig:spe_trajectory} places these record cells, with
the post hoc oracle of Section~\ref{sec:spe_discussion}, on one plane of accuracy against prompt length
for both Gemma scales.

\begin{figure}[htbp]
  \centering
  \includegraphics[width=\textwidth]{results/thesis_visuals/fig10_spe_story.pdf}
  \caption{Accuracy against mean prompt tokens for Gemma~4 E2B and E4B on 200 LoCoMo questions, with
  corrected prefixes and the repaired selector. Whiskers are Wilson 95\% intervals. Hexagons
  show the post hoc oracle selecting between v3 and v4 using judged labels; they bound those
  choices rather than represent an achievable policy. S2-b240 and S2-b370 mark SPE-2 at
  budgets $B=240$ and $B=370$. SPE-1 is omitted for legibility and
  reported in Tables~\ref{tab:spe_results} and~\ref{tab:scale_comparison}.}
  \label{fig:spe_trajectory}
\end{figure}

The two columns change extractor and reader together; the crossover in
Table~\ref{tab:ch6_crossover} associates most of the v3 difference with the graph producer
rather than reader scale.

\subsection{Why SPE retained only triples}
\label{sec:spe_why_triples}
Across both models and both budgets, none of the 800 pre-specified SPE-2 prompts (the registered
E2B study and its frozen E4B addendum) selected a source turn.\footnote{Only the
repaired record runs are counted here; the earlier selector is described in
Section~\ref{sec:spe_results}.} The exploratory SPE-1 variant selected none either. At
$B=370$, E2B retained 20.7 triples on average and E4B retained 18.9. SPE-2 was therefore a larger
triple-only retriever, not a selective expansion policy.

Three properties of the inputs explain this outcome: what source text costs, what an upgrade
adds, and how many triples were available.

\paragraph{Cost of source text.} Source text was much more expensive than a triple. Measured offline from the
logged prompts and graph caches, in the selector's own token units, the 4,136 triples that E2B
SPE-2 selected at $B=370$ cost 17.5 tokens each on average. Upgrading one of them would have cost
75.4 more tokens (median 68), because a quarter of these facts pointed to two or more retrieved
turns. The median ratio $r=\Delta C/C$ of upgrade cost to triple cost was 3.9 for E2B and 3.5 for
E4B, and $r$ was below two for only 10\% of E2B upgrades and 13\% of E4B upgrades
(\texttt{results/spe\_upgrade\_cost.json}).

\paragraph{What an upgrade adds.} At each step the selector compares what each option adds per token. A new
triple with similarity $u'$ to the query and cost $C'$ adds its full similarity, $u'/C'$. An
upgrade adds only the gap,
$\text{gap}(e)/\Delta C$, because the triple's similarity is already in the prompt. For a competing
triple of similar length, the upgrade wins only if
\begin{equation}
  \text{gap}(e) > r \cdot u' .
\end{equation}
The figure of about 0.26 below follows from this inequality under four stated assumptions:
\begin{enumerate}
  \item the cost ratio $r$ takes its median value of 3.9 (E2B); for an individual fact it can be
        smaller or larger;
  \item the competing triple is about as long as the triple being upgraded, so $C' \approx C$;
  \item both the triple and its source turn have non-negative similarity to the query, so the gap,
        which is the source turn's similarity minus the triple's, is below the source turn's
        similarity and therefore below 1;
  \item the greedy step compares the upgrade with the best remaining triple that still fits.
\end{enumerate}
Under these assumptions an upgrade can win only when every competing triple has a similarity below
$\text{gap}(e)/r < 1/3.9 \approx 0.26$, and in practice below a quarter of the actual gap. Adding a
fact directly in expanded form faces a similar bar: it costs the triple and the source text
together, about 5.3 times a triple, so under the same assumptions the competing triples would need a
similarity below about 0.19. The threshold is a consequence of the scoring rule, not a measured
value: the per-candidate similarities were not logged.

\paragraph{Number of triples.} An upgrade could also have won by default if the triples had run out
with budget left, but they never ran out. Including every candidate triple would take at least 727
tokens on average for E2B and 814 for E4B, about twice $B=370$. The selector logs confirm this for
every prompt: in none of the 800 SPE-2 prompts did it take every candidate triple, and it stopped only
after spending at least 220 of 240 and 348 of 370 tokens (\texttt{spe\_spent\_tokens} and
\texttt{candidate\_edges} in the run records). The logs show only that this default route never
opened. They do not show why each upgrade lost its comparison with a triple; that part of the
explanation rests on the scoring calculation above. Finally, the selector discards options whose
utility is not positive, so a fact whose triple was at least as similar to the query as its source
turn had no upgrade option at all.

The first expansion is the difficult step. Once one source turn is in the prompt, upgrading another fact
that points to the same turn costs only one token, but under the repaired rule that first expansion
never happened. The earlier version did reach it in 39 E2B prompts because it credited an upgrade
with the source turn's full similarity instead of the gap (Section~\ref{sec:spe_repair}). Its
condition was $u_{\text{source}} > r \cdot u'$, which is easier to meet because the source turn's
similarity is larger than the gap.

This explanation follows from the rule and the measured costs. The runs did not log the
utilities, so how often gaps were non-positive, and how close the best upgrades came to the bar,
remain unmeasured; establishing either requires logged utilities or an ablation of the selector.

The representation-specific scoring also added query-time work. Mean retrieval time on the workstation was
about 7.1 seconds per question for both SPE-2 budgets, compared with 0.21 seconds for the SPE-1
cells, 0.19 seconds for E2B GraphRAG v4 and 0.24 seconds for E2B Vector RAG. This implementation of
the selection rule therefore carried a substantial latency cost.

\section{Discussion}
\label{sec:spe_discussion}

In both Gemma configurations, the selector did not recover the value of the omitted dialogue text.
That conclusion concerns the selection rule, not whether a useful middle point existed; a post hoc
oracle bound separates the two. For each question it keeps the v4 answer when v4 was judged correct
and v3 was not, and the cheaper v3 answer otherwise. For E2B this reaches 35.0\% at a mean of 560.3
prompt tokens, 31.1\% fewer than v4, which would satisfy Criterion~1; for E4B it reaches 33.5\% at
555.6 tokens, 29.7\% fewer than v4. Expanding the same share of questions at random would give 23.1\%
and 27.8\%. The oracle reads the judged labels, so it is an upper bound on question-level selection
between two existing configurations rather than an achievable policy. Its value is to show where the
headroom lay: source restoration turned 27 E2B answers and 12 E4B answers from wrong to right while
reversing 4 and 2, so a selector had to find a small minority of questions. The tested rule never
spent its budget on source text, so it could not find them.

For these graphs and this dialogue QA task, the tested greedy policies did not identify a useful
middle point between compact triples and source text. Text-attributed graph nodes might perform
differently, although adding longer text to each node would narrow the practical distinction from
chunk-based retrieval while retaining the graph-construction cost measured in
Section~\ref{sec:ch5_graph_boundary}.

SEEM and LeanMem also return from a compact memory unit to its source text
(Section~\ref{sec:ch2_graphs}): SEEM for every
retrieved event frame, within twice the initial retrieval budget~\citep{seem}, and LeanMem
only when an LLM planner judges that a query needs detailed evidence~\citep{leanmem}. Both
are evaluated with server-scale backbones. SPE differs in where the decision is made: it
weighs the two representations of each retrieved fact against each other under a shared
token budget. That per-fact comparison is what failed here. A query-level decision of the
kind LeanMem makes, for which the oracle bound above shows headroom, was not tested.

\section{Limitations}
\label{sec:spe_limitations}

The limits of the LoCoMo sample, the automatic judge and graph extraction stated in
Chapters~\ref{ch:memory_evaluation} and~\ref{ch:graph_mechanisms} apply here, and SPE ran on the
workstation only. Three limits are specific to SPE:
\begin{enumerate}
  \item \textbf{Selection policies}: SPE tested one heuristic and one representation-specific cosine
        score, both ranked by utility per token. That normalization was never ablated. Cosine
        similarity is not a ratio-scale measure of usefulness, so dividing it by token cost favors the
        shorter representation (Section~\ref{sec:spe_why_triples}). Calibrated utilities, learned
        policies, cross-encoder reranking, a budget share reserved for source turns and a query-level
        expansion decision were not tested.
  \item \textbf{Scope}: All SPE runs used Gemma~4, E2B for the registered study and E4B for the frozen
        addendum, at budgets of 240 and 370 tokens. The negative result does not establish cross-family
        generality; testing another family now would be a new replication, not part of the frozen
        evaluation.
  \item \textbf{Post hoc analyses}: The oracle bound reads judged labels, and the reader--extractor
        crossover was added after the scale comparison exposed a confound. Both are exploratory
        diagnostics, not deployable policies or confirmatory tests.
\end{enumerate}

\section{Summary}
\label{sec:spe_summary}

Selective Provenance Expansion did not provide a reliable middle point between triple-only retrieval
and source dialogue for the tested Gemma configurations. On E2B, Criterion~1 failed, and Criterion~2
cleared its numerical threshold without a statistically distinguishable gain; E4B also stayed below
the source-bearing baselines. None of the 800 pre-specified SPE-2 prompts selected a source turn: the
density rule spent the budget on additional triples. A post hoc oracle shows question-level headroom
between v3 and v4, so the negative result applies to the tested selection rule, not to selective
expansion as a whole.
